% Luc Destombe --- licence CC BY-NC-SA 4.0
% https://creativecommons.org/licenses/by-nc-sa/4.0/deed.fr
% Fichier autonome : compiler avec pdflatex, deux fois (signets, nombre de pages).
\documentclass[11pt,a4paper]{article}
\makeatletter
\newif\ifeleve
\elevefalse
\elevetrue

\newif\iflycee@palatino
\lycee@palatinotrue

\RequirePackage[utf8]{inputenc}
\RequirePackage[T1]{fontenc}
\RequirePackage[french]{babel}
\RequirePackage{etoolbox}

\newcommand{\ShorthandsOff}{\@ifundefined{shorthandoff}{}{\shorthandoff{;:!?}}}
\newcommand{\ShorthandsOn}{\@ifundefined{shorthandon}{}{\shorthandon{;:!?}}}
\ShorthandsOff
\AtBeginDocument{\ShorthandsOn}
\AtBeginEnvironment{tikzpicture}{\ShorthandsOff}

\RequirePackage{geometry}
\RequirePackage{fancyhdr}
\RequirePackage{lastpage}
\RequirePackage{multicol}
\RequirePackage{array}
\RequirePackage{enumitem}
\RequirePackage{xcolor}
\RequirePackage{needspace}

\RequirePackage{amsmath,amssymb}
\RequirePackage{mathpazo}
\DeclareMathSizes{10.95}{10.95}{8}{5.5}
\begingroup\catcode`\_=\active \catcode`\^=\active
\gdef_#1{\sb{\scriptscriptstyle #1}}
\gdef^#1{\sp{\scriptscriptstyle\lycee@moinsepais #1}}
\endgroup
\newcommand\lycee@moins{\mathbin{%
  \setbox\z@\hbox{$\m@th\scriptscriptstyle\lycee@moinsorig$}%
  \dimen@=.55\wd\z@ \advance\dimen@-.5pt
  \hbox to.75\wd\z@{\hss$\m@th\scriptscriptstyle\vcenter{\hbox{%
    \tikz\draw[line width=.5pt,line cap=round](0,0)--(\the\dimen@,0);}}$\hss}}}
\begingroup\lccode`\~=`\- \lowercase{\endgroup\let~}\lycee@moins
\newcommand\lycee@moinsepais{\mathcode`\-="8000 }
\AtBeginDocument{\mathchardef\lycee@moinsorig=\mathcode`\-\relax}
\AtBeginDocument{\catcode`\_=12 \mathcode`\_="8000
  \catcode`\^=12 \mathcode`\^="8000 }
\newcommand\lycee@exposants{%
  \fontdimen13\textfont2=.48\fontdimen6\textfont2
  \fontdimen14\textfont2=.48\fontdimen6\textfont2
  \fontdimen15\textfont2=.48\fontdimen6\textfont2
  \fontdimen13\scriptfont2=.48\fontdimen6\scriptfont2
  \fontdimen14\scriptfont2=.48\fontdimen6\scriptfont2
  \fontdimen15\scriptfont2=.48\fontdimen6\scriptfont2}
\every@math@size\expandafter{\the\every@math@size\lycee@exposants}
\newlength{\retraitcalculs}
\setlength{\retraitcalculs}{2em}

\RequirePackage{tikz}
\usetikzlibrary{arrows.meta,calc,babel,shapes.misc}
\RequirePackage[most]{tcolorbox}
\tcbuselibrary{skins,breakable}

\definecolor{coulDef}{HTML}{1F4E79}
\definecolor{coulProp}{HTML}{7030A0}
\definecolor{coulRem}{HTML}{C55A11}
\definecolor{coulEx}{HTML}{2E7D32}
\definecolor{coulExo}{HTML}{B03A2E}
\definecolor{coulPart}{HTML}{1F4E79}
\definecolor{coulMeth}{HTML}{1B6E4F}
\definecolor{coulCourbe}{HTML}{0B5ED7}
\definecolor{coulCourbeB}{HTML}{C00000}
\definecolor{coulCourbeC}{HTML}{2E7D32}
\definecolor{coulGrille}{HTML}{8C8C8C}
\definecolor{coulRep}{HTML}{1B6E4F}
\definecolor{coulBar}{HTML}{2E7D32}

\newcommand{\R}{\mathbb{R}}
\newcommand{\Rs}{\mathbb{R}^{*}}
\renewcommand{\le}{\leqslant}
\renewcommand{\ge}{\geqslant}
\renewcommand{\approx}{\simeq}
\newcommand{\vect}[1]{\overrightarrow{#1}}
\newcommand{\norme}[1]{\left\lVert #1 \right\rVert}
\newcommand{\intervalle}[2]{\left[#1\,;#2\right]}
\RequirePackage{cancel}
\renewcommand{\CancelColor}{\color{coulExo}}
\newcommand{\fois}[1]{\times{\color{coulExo}#1}}
\newcommand{\barre}[1]{\cancel{#1}}

\tikzset{grille/.style={coulGrille,line width=0.4pt}}
\newcommand{\quadrillage}[4]{\draw[grille,step=1] (#1,#2) grid (#3,#4);}
\tikzset{
  croix/.style={cross out,draw,line width=0.6pt,inner sep=0pt,minimum size=4pt},
  croixdroite/.style={croix,rotate=45,line width=0.8pt},
}
\newcommand{\pt}[3]{\path (#1) node[croix]{} node[#2,font=\small,inner sep=2.5pt]{$#3$};}

\newcommand{\lycee@repere}[4]{%
  \draw[grille,step=1] (#1,#2) grid (#3,#4);
  \draw[-{Stealth[length=2mm]},thick] (#1,0)--({#3+0.4},0) node[below left=-1pt]{$x$};
  \draw[-{Stealth[length=2mm]},thick] (0,#2)--(0,{#4+0.4}) node[below left=-1pt]{$y$};
  \node[below left,font=\scriptsize,inner sep=1.5pt] at (0,0) {$0$};}
\newcommand{\lycee@gradx}[1]{\foreach \k/\l in {#1}{%
  \draw (\k,0.07)--(\k,-0.07) node[below,font=\scriptsize,inner sep=1.5pt]{$\l$};}}
\newcommand{\lycee@grady}[1]{\foreach \k/\l in {#1}{%
  \draw (0.07,\k)--(-0.07,\k) node[left,font=\scriptsize,inner sep=1.5pt]{$\l$};}}
\newcommand{\lycee@intff}[2]{\left[#1\,;#2\right]}
\newcommand{\lycee@intfo}[2]{\left[#1\,;#2\right[}
\newcommand{\lycee@intof}[2]{\left]#1\,;#2\right]}
\newcommand{\lycee@intoo}[2]{\left]#1\,;#2\right[}
\newcommand{\lycee@ens}[1]{\left\{#1\right\}}
\AtBeginDocument{%
  \providecommand{\repere}{\lycee@repere}%
  \providecommand{\gradx}{\lycee@gradx}%
  \providecommand{\grady}{\lycee@grady}%
  \providecommand{\Cf}{\mathcal{C}\sb{\scriptscriptstyle f}}%
  \providecommand{\Cg}{\mathcal{C}\sb{\scriptscriptstyle g}}%
  \providecommand{\intff}{\lycee@intff}%
  \providecommand{\intfo}{\lycee@intfo}%
  \providecommand{\intof}{\lycee@intof}%
  \providecommand{\intoo}{\lycee@intoo}%
  \providecommand{\ens}{\lycee@ens}}

\newlist{questions}{enumerate}{3}
\setlist[questions,1]{label=\arabic*\textdegree),leftmargin=*,itemsep=2pt,topsep=3pt}
\setlist[questions,2]{label=\alph*),leftmargin=*,itemsep=1pt,topsep=2pt}
\setlist[questions,3]{label=\roman*),leftmargin=*,itemsep=1pt,topsep=1pt}
\setlist[itemize]{label=\textbullet,leftmargin=*,itemsep=2pt,topsep=3pt}

\newcommand{\besoinplace}[1]{\par\penalty\@M\begingroup
  \dimen@=#1\relax\dimen@ii=\pagegoal\advance\dimen@ii-\pagetotal
  \ifdim\dimen@>\dimen@ii\ifdim\dimen@ii>\z@\vfil\fi\break\fi\endgroup}

\newcommand{\lycee@pied}{\thepage/\pageref{LastPage}}
\newcommand{\entete}[2]{%
  \pagestyle{fancy}\fancyhf{}%
  \fancyhead[L]{\footnotesize\color{coulPart}\textbf{#1}}%
  \fancyhead[R]{\footnotesize\color{coulPart}\textbf{#2}}%
  \fancyfoot[C]{\footnotesize\lycee@pied}%
  \renewcommand{\headrulewidth}{0.4pt}%
  \renewcommand{\headrule}{\hbox to\headwidth{%
    \color{coulPart!40}\leaders\hrule height \headrulewidth\hfill}}}

\RequirePackage{ccicons}
\newcommand{\lycee@licence}{{\scriptsize\color{gray}%
  \copyright\ Luc Destombe --- \ccbyncsa\ CC BY-NC-SA 4.0}}
\newif\iflycee@licence \lycee@licencetrue
\newcommand{\sanslicence}{\lycee@licencefalse}
\AtBeginDocument{\iflycee@licence\AddToHookNext{shipout/foreground}{%
  \put(0,0){\rlap{\kern\dimexpr1in+\hoffset+\oddsidemargin+\textwidth\relax
    \llap{\raisebox{-\dimexpr1in+\voffset+\topmargin+\headheight+\headsep
      +\textheight+\footskip\relax}{\lycee@licence}}}}}\fi}

\newcommand{\formulesserrees}{%
  \setlength{\abovedisplayskip}{3pt}\setlength{\belowdisplayskip}{3pt}%
  \setlength{\abovedisplayshortskip}{2pt}\setlength{\belowdisplayshortskip}{3pt}}

\setlength{\parindent}{0pt}

\RequirePackage[implicit=false,hidelinks,pdfencoding=auto,psdextra,bookmarksopen,
  bookmarksopenlevel=1,pdfcreator={},pdfproducer={}]{hyperref}
\RequirePackage{bookmark}
\hypersetup{pdfauthor={Luc Destombe},pdfkeywords={CC BY-NC-SA 4.0}}
\newcounter{lycee@signet}
\newcommand{\signet}[3][]{\stepcounter{lycee@signet}%
  \edef\lycee@dest{\if\relax\detokenize{#1}\relax
    signet.\arabic{lycee@signet}\else#1\fi}%
  \hypertarget{\lycee@dest}{}\bookmark[level=#2,dest=\lycee@dest]{#3}}
\pdfstringdefDisableCommands{%
  \def\R{R}\def\vect#1{#1}\def\Cf{Cf}\def\Cg{Cg}\def\fois#1{×#1}%
  \def\barre#1{#1}\def\intervalle#1#2{[#1;#2]}\def\intff#1#2{[#1;#2]}%
  \def\textsuperscript#1{#1}\def\dfrac#1#2{#1/#2}\def\frac#1#2{#1/#2}%
  \def\vec#1{#1⃗}}%
\geometry{top=2cm,bottom=1cm,left=1cm,right=1cm,headheight=14pt,footskip=0.6cm}

\RequirePackage{pgfplots}
\pgfplotsset{compat=1.18}
\RequirePackage{tkz-tab}

\newcounter{partiecours}
\renewcommand{\thepartiecours}{\Roman{partiecours}}
\newcommand{\partiecours}[1]{%
  \refstepcounter{partiecours}%
  \Needspace*{8\baselineskip}%
  \bigskip
  \noindent\begin{tcolorbox}[enhanced,colback=coulPart,colframe=coulPart,
      boxrule=0pt,arc=2pt,left=8pt,right=8pt,top=4pt,bottom=4pt,
      before skip=6pt,after skip=10pt]
    \color{white}\bfseries\large \thepartiecours{} --- #1%
    \signet[partie-\arabic{partiecours}]{1}{\thepartiecours{} --- #1}
  \end{tcolorbox}%
  \addcontentsline{toc}{section}{\thepartiecours{} --- #1}%
}

\newtcolorbox{boitecours}[3][]{%
  enhanced,breakable,
  colback=white,colframe=#2,coltitle=white,
  fonttitle=\bfseries,
  attach boxed title to top left={xshift=8pt,yshift=-\tcboxedtitleheight/2},
  boxed title style={colback=#2,arc=2pt,boxrule=0pt},
  boxrule=0.9pt,arc=2pt,
  left=8pt,right=8pt,top=8pt,bottom=6pt,
  title={#3},#1}

\newcounter{methode}
\newenvironment{definitionbox}[1][Définition]%
  {\begin{boitecours}[phantom={\signet{2}{#1}}]{coulDef}{#1}}{\end{boitecours}}
\newenvironment{proprietebox}[1][Propriété]%
  {\begin{boitecours}[phantom={\signet{2}{#1}}]{coulProp}{#1}}{\end{boitecours}}
\newenvironment{methodebox}[1][Méthode]{\stepcounter{methode}%
  \begin{boitecours}[phantom={\signet[methode-\arabic{methode}]{2}{#1}}]{coulMeth}{#1}}%
  {\end{boitecours}}

\newcommand{\etape}[1]{\par\smallskip\textbf{\color{coulMeth}#1}\par\nobreak\smallskip}
\newcommand{\enonce}[1]{\emph{#1}\par\smallskip}
\newcommand{\reponse}[1]{\par\smallskip
  {\footnotesize\itshape\color{black!55}Réponses : #1}}
\newcommand{\demo}[1]{\textbf{\color{coulMeth}Démonstration.} #1}
\newcommand{\afaire}[1]{%
  \par\smallskip
  \noindent{\small\color{coulExo}$\blacktriangleright$\ \textbf{À faire :}
    fiche d'exercices du chapitre, #1.}\par\smallskip}

\newenvironment{calculs}{\setlength{\jot}{5pt}\@fleqntrue\@mathmargin\retraitcalculs
  \setlength{\abovedisplayskip}{4pt}\setlength{\belowdisplayskip}{4pt}%
  \setlength{\abovedisplayshortskip}{4pt}\setlength{\belowdisplayshortskip}{4pt}%
  \csname alignat*\endcsname{2}}%
  {\csname endalignat*\endcsname}
\newcommand{\rem}[1]{\text{\small\itshape\color{black!60}#1}}
\newcommand{\explique}[2]{\par\smallskip\noindent
  \begin{minipage}[t]{0.57\linewidth}\raggedright #1\end{minipage}\hfill
  \begin{minipage}[t]{0.40\linewidth}\raggedright\leavevmode
    {\small\itshape\color{black!60}#2}\end{minipage}\par}
\newcommand{\cas}[1]{\par\smallskip\noindent\textbf{\color{coulExo}#1}\nobreak}
\newcommand{\calc}[1]{\par\smallskip\textbf{\color{coulExo}#1}\par\nobreak}

\newtcolorbox{remarquebox}[1][Remarque]{%
  enhanced,breakable,colback=white,colframe=coulRem,
  boxrule=0pt,leftrule=2.5pt,arc=0pt,outer arc=0pt,
  left=8pt,right=6pt,top=5pt,bottom=5pt,
  before upper={\textbf{\color{coulRem}#1}\par\smallskip}}

\newtcolorbox{exemplebox}[1][Exemple]{%
  enhanced,breakable,colback=white,colframe=coulEx,
  boxrule=0pt,leftrule=2.5pt,arc=0pt,outer arc=0pt,
  left=8pt,right=6pt,top=5pt,bottom=5pt,
  before upper={\textbf{\color{coulEx}#1}\par\smallskip}}

\pgfplotsset{
  repere/.style={
    axis lines=middle,
    axis line style={-{Stealth[length=2.4mm]},thick},
    every axis x label/.style={at={(ticklabel* cs:1.0)},anchor=west},
    every axis y label/.style={at={(ticklabel* cs:1.0)},anchor=south},
    label style={font=\footnotesize},
    tick label style={font=\scriptsize},
    grid=both,
    major grid style={grille},
    minor tick num=0,
    tick align=inside,
    clip mode=individual,
  },
}
\tikzset{
  courbe/.style={coulCourbe,line width=1.1pt,smooth},
  courbeB/.style={coulCourbeB,line width=1.1pt,smooth},
  courbeC/.style={coulCourbeC,line width=1.1pt,smooth},
}

\newlist{etapes}{enumerate}{1}
\setlist[etapes]{label=\textbf{\arabic*.},leftmargin=*,itemsep=1pt,topsep=2pt}

\newlist{expressions}{enumerate}{1}
\setlist[expressions]{label={\Alph*\ $=$},leftmargin=*,itemsep=3pt,topsep=3pt}

\newcounter{exo}
\renewcommand{\theexo}{\ifnum\value{exo}<10 0\fi\arabic{exo}}
\newtcolorbox{exercice}[1][]{%
  enhanced,breakable,
  colback=white,colframe=coulExo,
  boxrule=0.9pt,arc=2pt,
  left=8pt,right=8pt,top=8pt,bottom=6pt,
  coltitle=white,fonttitle=\bfseries,
  attach boxed title to top left={xshift=8pt,yshift=-\tcboxedtitleheight/2},
  boxed title style={colback=coulExo,arc=2pt,boxrule=0pt},
  phantom={\signet{2}{Exercice \arabic{exo}}},
  title={Exercice \theexo\ifx\relax#1\relax\else{} --- #1\fi}}
\newenvironment{exo}[1][\relax]{\refstepcounter{exo}\begin{exercice}[#1]}{\end{exercice}}

  \RequirePackage{comment}
  \excludecomment{correction}
  \renewcommand{\reponse}[1]{}
  \newcommand{\autableau}{\hfill{\footnotesize\itshape\color{black!45}(au tableau)}}
  \newcommand{\etapeexemples}{%
    \par\smallskip\textbf{\color{coulMeth}Exemples}\autableau\par\nobreak\smallskip}
  \newcommand{\etapetableau}[1]{%
    \par\smallskip\textbf{\color{coulMeth}#1}\autableau\par\nobreak\smallskip}
  \newtcolorbox{exempletableau}[1][Exemple]{%
    enhanced,breakable,colback=white,colframe=coulEx,
    boxrule=0pt,leftrule=2.5pt,arc=0pt,outer arc=0pt,
    left=8pt,right=6pt,top=4pt,bottom=4pt,
    before skip=5pt,after skip=5pt,
    before upper={\textbf{\color{coulEx}#1}\autableau\par\smallskip}}

\RequirePackage{comment}
\excludecomment{correctionavj}

\newcommand{\coursEntete}{}
\newcommand{\coursNiveau}{}
\newcommand{\coursTitre}{}
\newcommand{\coursSurTitre}{}
\newcommand{\coursSousTitre}{}

\newcommand{\enteteCours}[1]{\renewcommand{\coursEntete}{#1}}
\newcommand{\chapitrecours}[2]{\enteteCours{Chapitre #1 --- #2}}
\newcommand{\niveaucours}[1]{\renewcommand{\coursNiveau}{#1}}
\newcommand{\titrecours}[3][]{%
  \renewcommand{\coursTitre}{#2}%
  \renewcommand{\coursSurTitre}{#1}%
  \renewcommand{\coursSousTitre}{#3}}

\pagestyle{fancy}
\fancyhf{}
\fancyhead[L]{\footnotesize\color{coulPart}\textbf{\coursEntete}}
\fancyhead[R]{\footnotesize\color{coulPart}\textbf{\coursNiveau}}
\fancyfoot[C]{\footnotesize\thepage}
\renewcommand{\headrulewidth}{0.4pt}
\renewcommand{\headrule}{\hbox to\headwidth{\color{coulPart!40}\leaders\hrule height \headrulewidth\hfill}}

\setlength{\parskip}{2pt}

\AtBeginDocument{%
  \begin{center}
    \begin{tcolorbox}[enhanced,width=0.72\linewidth,colback=white,colframe=coulPart,
        boxrule=1.6pt,arc=3pt,halign=center,top=10pt,bottom=10pt]
      \ifx\coursSurTitre\empty
        {\Huge\bfseries\color{coulPart} \coursTitre}\\[4pt]
      \else
        {\Huge\bfseries\color{coulPart} \coursTitre}\\[3pt]
        {\Large\bfseries\color{coulPart} \coursSurTitre}\\[5pt]
      \fi
      {\normalsize \coursSousTitre}%
      \\[2pt]{\small\itshape Version élève}
    \end{tcolorbox}
  \end{center}%
}
\makeatother

\chapitrecours{3}{Calcul littéral et second degré}
\niveaucours{1\textsuperscript{re} mathématiques spécifiques}
\titrecours{CALCUL LITTÉRAL ET SECOND DEGRÉ}{Cours --- Première, mathématiques spécifiques}

\begin{document}

\begin{remarquebox}[Comment utiliser ce cours]
Chaque partie commence par ce qu'il faut \textbf{savoir} (définitions, propriétés), puis
vient ce qu'il faut \textbf{savoir faire} : une méthode, avec des
\textbf{exemples} faits en classe, puis des exercices
\og\textbf{À vous de jouer}\fg{} à chercher seul. Tout le chapitre se fait
\textbf{sans calculatrice}.

\smallskip
Sur cette feuille, seuls les \textbf{énoncés} des exemples figurent : la résolution,
signalée par la mention \emph{(au tableau)}, est à rédiger dans le cahier.

\end{remarquebox}

\partiecours{Développer}

\begin{definitionbox}[Définition 1 --- Développer]
\begin{minipage}[t]{0.46\linewidth}
\vspace{2pt}
\textbf{Développer}, c'est transformer un \textbf{produit} en une \textbf{somme}.

\smallskip
Le chemin inverse s'appelle \textbf{factoriser}.
\end{minipage}\hfill
\begin{minipage}[t]{0.50\linewidth}
\vspace{0pt}
\centering
\begin{tikzpicture}[>={Stealth[length=3mm]}]
  \node (g) at (0,0) {$3x(x+4)$};
  \node (d) at (3.4,0) {$3x^{2}+12x$};
  \draw[->,line width=1pt,coulEx] (g.north east) to[bend left=32]
        node[above,font=\scriptsize\bfseries,coulEx]{DÉVELOPPER} (d.north west);
  \draw[->,line width=1pt,coulExo] (d.south west) to[bend left=32]
        node[below,font=\scriptsize\bfseries,coulExo]{FACTORISER} (g.south east);
\end{tikzpicture}
\end{minipage}
\end{definitionbox}

\begin{proprietebox}[Propriété 1 --- Distributivité]
Pour tous nombres $k$, $a$, $b$, $c$ et $d$ :
\[
  k(a+b)=ka+kb \qquad\qquad (a+b)(c+d)=ac+ad+bc+bd
\]
Dans $(a+b)(c+d)$, chaque terme de la première parenthèse multiplie chaque terme de
la seconde : quatre produits.
\end{proprietebox}

\begin{methodebox}[Méthode 1 --- Développer et réduire]
\etapeexemples
Développer et réduire :
\begin{questions}[label=\alph*)]
  \item $A=3(2x-5)$
  \item $B=-2x(x+4)$
  \item $C=(x+3)(2x-1)$
\end{questions}

\etape{À vous de jouer}
Développer et réduire :
\[
  D=4(3x+2) \qquad E=-5(x-3) \qquad F=-3x(x+2) \qquad G=x(5-x)
\]
\[
  H=(x+4)(x-2) \qquad I=(2x-3)(x+5)
\]

\end{methodebox}

\afaire{exercices 1 à 4}

\partiecours{Identités remarquables}

\begin{proprietebox}[Propriété 2 --- Identités remarquables]
Pour tous nombres $a$ et $b$ :
\[
  (a+b)^{2}=a^{2}+2ab+b^{2} \qquad (a-b)^{2}=a^{2}-2ab+b^{2} \qquad
  (a+b)(a-b)=a^{2}-b^{2}
\]
Attention : $(x+5)^{2}$ n'est pas égal à $x^{2}+25$ ; il manque le \textbf{double produit}
$2ab$.
\end{proprietebox}

\begin{methodebox}[Méthode 2 --- Développer avec une identité remarquable]
\etapeexemples
Développer et réduire :
\begin{questions}[label=\alph*)]
  \item $A=(x+5)^{2}$
  \item $B=(3x-2)^{2}$
  \item $C=(x+7)(x-7)$
\end{questions}

\etape{À vous de jouer}
Développer et réduire :
\[
  D=(x+3)^{2} \qquad E=(2x+5)^{2} \qquad F=(x-8)^{2} \qquad G=(4x-1)^{2}
\]
\[
  H=(x+10)(x-10) \qquad I=(2x-9)(2x+9)
\]

\end{methodebox}

\afaire{exercices 5 à 8}

\partiecours{Équations du premier degré}

\begin{definitionbox}[Définition 2 --- Résoudre une équation]
\textbf{Résoudre une équation} d'inconnue $x$, c'est trouver toutes les valeurs de $x$
qui rendent l'égalité vraie : ce sont les \textbf{solutions}. On les écrit entre
accolades : $S=\ens{4}$.
\end{definitionbox}

\begin{proprietebox}[Propriété 3 --- Transformer une équation]
On obtient une équation qui a les mêmes solutions quand on \textbf{ajoute} (ou
\textbf{retranche}) un même nombre aux deux membres, ou quand on les
\textbf{multiplie} (ou les \textbf{divise}) par un même nombre non nul.
\end{proprietebox}

\begin{methodebox}[Méthode 3 --- Résoudre une équation du premier degré]
\etapeexemples
\begin{questions}[label=\alph*)]
  \item Résoudre $7x-3=4x+9$.
  \item Résoudre $2(x+1)=5x-7$.
  \item Résoudre $\dfrac{20}{x}=4$.
  \item Une salle d'escalade propose deux tarifs : $8$~€ la séance, ou un abonnement
        de $30$~€ puis $5$~€ la séance. Pour combien de séances les deux tarifs
        coûtent-ils autant ?
\end{questions}

\etape{À vous de jouer}
\begin{questions}[label=\alph*)]
  \item Résoudre $6x+5=2x+17$, puis $9x-4=-x+6$.
  \item Résoudre $3(x-2)=x+8$, puis $-2(x-4)=x-1$.
  \item Résoudre $\dfrac{36}{x}=9$, puis $\dfrac{15}{x}=-3$.
  \item Une location de trottinette coûte $4$~€ par jour, ou $18$~€ d'inscription
        puis $1$~€ par jour. Pour combien de jours les deux formules coûtent-elles
        autant ?
\end{questions}

\end{methodebox}

\afaire{exercices 9 à 12}

\partiecours{Équations $x^{2}=a$}

\begin{proprietebox}[Propriété 4 --- Équation $x^{2}=a$]
\begin{minipage}[c]{0.58\linewidth}\raggedright
\begin{itemize}
  \item Si $a>0$ : deux solutions, $-\sqrt{a}$ et $\sqrt{a}$.
  \item Si $a=0$ : une seule solution, $0$.
  \item Si $a<0$ : aucune solution, car un carré n'est jamais négatif.
\end{itemize}
\smallskip
Sur le graphique : la droite horizontale $y=a$ coupe la courbe de $x\mapsto x^{2}$ en
deux points si $a>0$, en aucun si $a<0$.
\end{minipage}\hfill
\begin{minipage}[c]{0.38\linewidth}\centering
\begin{tikzpicture}[x=0.5cm,y=0.5cm]
  \repere{-3}{-2}{3}{5}
  \grady{1,2,3,4}
  \draw[coulCourbe,line width=1.1pt,domain=-2.2:2.2,samples=60] plot (\x,{\x*\x});
  \draw[red,dashed,line width=0.9pt] (-3,4)--(3,4);
  \node[red,above right,font=\footnotesize,inner sep=1pt] at (-3,4) {$y=a$};
  \draw[red,dotted,thick] (-2,4)--(-2,0) (2,4)--(2,0);
  \path[red] (-2,4) node[croix]{} (2,4) node[croix]{};
  \node[red,below,font=\footnotesize,inner sep=2pt] at (-2,0) {$-\sqrt{a}$};
  \node[red,below,font=\footnotesize,inner sep=2pt] at (2,0) {$\sqrt{a}$};
  \draw[black!60,dashed,line width=0.9pt] (-3,-1)--(3,-1);
  \node[black!60,font=\footnotesize,fill=white,inner sep=1pt] at (1.9,-1.6) {$a<0$};
\end{tikzpicture}
\end{minipage}
\end{proprietebox}

\begin{methodebox}[Méthode 4 --- Résoudre une équation $x^{2}=a$]
\etapeexemples
\begin{questions}[label=\alph*)]
  \item Résoudre $x^{2}=36$.
  \item Résoudre $x^{2}=11$.
  \item Résoudre $x^{2}=-9$.
  \item Résoudre $2x^{2}-50=0$.
  \item Un jardin carré a une aire de $64~\text{m}^{2}$. Quelle est la longueur de son
        côté ?
\end{questions}

\etape{À vous de jouer}
\begin{questions}[label=\alph*)]
  \item Résoudre $x^{2}=49$, puis $x^{2}=100$.
  \item Résoudre $x^{2}=7$, puis $x^{2}=15$.
  \item Résoudre $x^{2}=-16$.
  \item Résoudre $3x^{2}-12=0$, puis $x^{2}+5=30$.
  \item Une nappe carrée a une aire de $4~\text{m}^{2}$. Quelle est la longueur de son
        côté ?
\end{questions}

\end{methodebox}

\afaire{exercices 13 à 15}

\partiecours{Équations produit}

\begin{proprietebox}[Propriété 5 --- Produit nul]
Un produit est nul quand l'un de ses facteurs est nul : si $A\times B=0$, alors $A=0$
ou $B=0$.
\end{proprietebox}

\begin{methodebox}[Méthode 5 --- Résoudre une équation produit]
\etapeexemples
Résoudre :
\begin{questions}[label=\alph*)]
  \item $(x-4)(x+6)=0$
  \item $(3x-6)(-2x+10)=0$
  \item $x^{2}-5x=0$
\end{questions}

\etape{À vous de jouer}
Résoudre :
\[
  (x-2)(x+9)=0 \qquad (x+1)(x+8)=0 \qquad (2x+8)(-3x+3)=0
\]
\[
  (5x-10)(4x+12)=0 \qquad x^{2}+3x=0 \qquad 2x^{2}-8x=0
\]

\end{methodebox}

\afaire{exercices 16 à 18}

\partiecours{Fonction polynôme du second degré}

\begin{definitionbox}[Définition 3 --- Fonction polynôme du second degré]
Une \textbf{fonction polynôme du second degré} est une fonction $f$ de la forme
\[
  f(x)=ax^{2}+bx+c \qquad\text{avec } a\ne 0 .
\]
Les nombres $a$, $b$ et $c$ sont ses \textbf{coefficients}. Sa courbe s'appelle une
\textbf{parabole}.
\end{definitionbox}

\begin{remarquebox}
Sans terme en $x^{2}$, il reste $bx+c$ : c'est une fonction affine, dont la courbe est
une droite.
\end{remarquebox}

\begin{methodebox}[Méthode 6 --- Reconnaître une fonction du second degré, calculer une image]
\etapeexemples
\begin{questions}[label=\alph*)]
  \item Pour chaque fonction $f$, $g$ et $h$, dire si elle est du second degré et, si
        oui, donner $a$, $b$ et $c$ :
        $f(x)=3x^{2}-2x+1$ \; ; \; $g(x)=5-x^{2}$ \; ; \; $h(x)=4x+7$.
  \item Calculer $f(-2)$.
  \item Les points $A(1\,;2)$ et $B(-1\,;4)$ sont-ils sur la courbe de $f$ ?
\end{questions}

\etape{À vous de jouer}
\begin{questions}[label=\alph*)]
  \item Pour chaque fonction $u$, $v$ et $w$, dire si elle est du second degré et, si
        oui, donner $a$, $b$ et $c$ :
        $u(x)=-2x^{2}+x+4$ \; ; \; $v(x)=x^{2}-3x$ \; ; \; $w(x)=1-6x$.
  \item Calculer $u(-1)$.
  \item Les points $C(2\,;-2)$ et $D(-2\,;0)$ sont-ils sur la courbe de $u$ ?
\end{questions}

\end{methodebox}

\afaire{exercices 19 à 22}

\partiecours{Forme factorisée et racines}

\begin{definitionbox}[Définition 4 --- Forme factorisée, racines]
\begin{minipage}[c]{0.58\linewidth}\raggedright
Une fonction du second degré est écrite sous \textbf{forme factorisée} quand elle
s'écrit $f(x)=a(x-x_{1})(x-x_{2})$.

\smallskip
Les \textbf{racines} de $f$ sont les solutions de $f(x)=0$. Ce sont les abscisses des
points où la parabole coupe l'axe des abscisses.
\end{minipage}\hfill
\begin{minipage}[c]{0.38\linewidth}\centering
\begin{tikzpicture}[x=0.42cm,y=0.42cm]
  \repere{-3}{-5}{5}{4}
  \gradx{-2,-1,1,2,3,4} \grady{-4,-2,2}
  \draw[coulCourbe,line width=1.1pt,domain=-1.9:3.9,samples=60] plot (\x,{(\x+1)*(\x-3)});
  \path[red] (-1,0) node[croix]{} (3,0) node[croix]{};
  \node[red,above left,font=\footnotesize,inner sep=1pt] at (-1,0) {$x_{1}$};
  \node[red,above right,font=\footnotesize,inner sep=1pt] at (3,0) {$x_{2}$};
\end{tikzpicture}
\end{minipage}
\end{definitionbox}

\begin{proprietebox}[Propriété 6 --- Racines d'une forme factorisée]
Les racines de $f(x)=a(x-x_{1})(x-x_{2})$ sont $x_{1}$ et $x_{2}$ : chacune annule une
parenthèse. Dans $(x-4)$, la racine est $4$ ; dans $(x+2)$, elle est $-2$.
\end{proprietebox}

\begin{methodebox}[Méthode 7 --- Utiliser une forme factorisée]
\etapeexemples
Soit $f(x)=2x^{2}+4x-6$.
\begin{questions}[label=\alph*)]
  \item Vérifier que $f(x)=2(x-1)(x+3)$.
  \item En déduire les racines de $f$.
  \item Donner les points où la courbe de $f$ coupe l'axe des abscisses.
\end{questions}

\etape{À vous de jouer}
Mêmes questions pour :
\begin{questions}[label=\alph*)]
  \item $g(x)=-x^{2}+2x+8$, avec $g(x)=-(x-4)(x+2)$ ;
  \item $h(x)=3x^{2}-12$, avec $h(x)=3(x-2)(x+2)$.
\end{questions}

\end{methodebox}

\afaire{exercices 23 à 25}

\partiecours{Représentation graphique}

\begin{proprietebox}[Propriété 7 --- Allure de la parabole]
La parabole de $f(x)=ax^{2}+bx+c$ :
\begin{itemize}
  \item est tournée vers le \textbf{haut} si $a>0$, vers le \textbf{bas} si $a<0$ ;
  \item a un \textbf{axe de symétrie} vertical, qui passe par son point le plus bas
        (ou le plus haut) : le \textbf{sommet} $S$ ;
  \item coupe l'axe des ordonnées au point $(0\,;c)$, car $f(0)=c$.
\end{itemize}
\begin{center}
\begin{minipage}[t]{0.47\linewidth}\centering
{\small $a>0$ : $f(x)=x^{2}-2x-1$}\par\smallskip
\begin{tikzpicture}[x=0.42cm,y=0.42cm]
  \repere{-2}{-3}{4}{4}
  \gradx{-1,1,2,3} \grady{-2,-1,1,2,3}
  \draw[coulCourbe,line width=1.1pt,domain=-1.4:3.4,samples=60] plot (\x,{\x*\x-2*\x-1});
  \draw[red,dashed,thick] (1,-3)--(1,4);
  \path (1,-2) node[croix]{} node[below right,font=\footnotesize,inner sep=2pt]{$S$};
  \path (0,-1) node[croix]{};
\end{tikzpicture}
\end{minipage}\hfill
\begin{minipage}[t]{0.47\linewidth}\centering
{\small $a<0$ : $g(x)=-x^{2}+2x+2$}\par\smallskip
\begin{tikzpicture}[x=0.42cm,y=0.42cm]
  \repere{-2}{-3}{4}{4}
  \gradx{-1,1,2,3} \grady{-2,-1,1,2,3}
  \draw[coulCourbe,line width=1.1pt,domain=-1.2:3.2,samples=60] plot (\x,{-\x*\x+2*\x+2});
  \draw[red,dashed,thick] (1,-3)--(1,4);
  \path (1,3) node[croix]{} node[above right,font=\footnotesize,inner sep=2pt]{$S$};
  \path (0,2) node[croix]{};
\end{tikzpicture}
\end{minipage}
\end{center}
\end{proprietebox}

\begin{methodebox}[Méthode 8 --- Tracer une parabole]
\etapeexemples
Soit $f(x)=x^{2}-2x-3$.
\begin{questions}[label=\alph*)]
  \item Compléter le tableau de valeurs de $f$ pour les entiers de $-2$ à $4$.
  \item Tracer la courbe de $f$.
  \item Lire le sommet et l'axe de symétrie.
\end{questions}

\etape{À vous de jouer}
Mêmes questions pour $g(x)=-x^{2}+4x$, avec les entiers de $-1$ à $5$.

\end{methodebox}

\begin{methodebox}[Méthode 9 --- Associer une fonction à sa parabole]
\etapeexemples
\begin{minipage}[c]{0.58\linewidth}\raggedright
Associer chaque fonction à sa courbe :
\[
  f(x)=x^{2}+1 \qquad g(x)=-x^{2}+3 \qquad h(x)=x^{2}-2
\]
\end{minipage}\hfill
\begin{minipage}[c]{0.38\linewidth}\centering
\begin{tikzpicture}[x=0.42cm,y=0.42cm]
  \repere{-3}{-3}{3}{5}
  \gradx{-2,2} \grady{-2,2,4}
  \draw[coulCourbe,line width=1.1pt,domain=-2.6:2.6,samples=60] plot (\x,{\x*\x-2});
  \draw[coulCourbeB,line width=1.1pt,domain=-2.2:2.2,samples=60] plot (\x,{-\x*\x+3});
  \draw[coulCourbeC,line width=1.1pt,domain=-2:2,samples=60] plot (\x,{\x*\x+1});
  \node[coulCourbe,font=\footnotesize] at (2.9,3.6) {$\mathcal{C}_{1}$};
  \node[coulCourbeB,font=\footnotesize] at (2.6,-2.4) {$\mathcal{C}_{2}$};
  \node[coulCourbeC,font=\footnotesize] at (-1.2,4.6) {$\mathcal{C}_{3}$};
\end{tikzpicture}
\end{minipage}

\etape{À vous de jouer}
\begin{minipage}[c]{0.58\linewidth}\raggedright
Associer chaque fonction à sa courbe :
\[
  p(x)=-x^{2}-1 \qquad q(x)=x^{2}-3 \qquad r(x)=-x^{2}+2
\]

\end{minipage}\hfill
\begin{minipage}[c]{0.38\linewidth}\centering
\begin{tikzpicture}[x=0.42cm,y=0.42cm]
  \repere{-3}{-5}{3}{4}
  \gradx{-2,2} \grady{-4,-2,2}
  \draw[coulCourbe,line width=1.1pt,domain=-2.6:2.6,samples=60] plot (\x,{-\x*\x+2});
  \draw[coulCourbeB,line width=1.1pt,domain=-2.6:2.6,samples=60] plot (\x,{\x*\x-3});
  \draw[coulCourbeC,line width=1.1pt,domain=-2:2,samples=60] plot (\x,{-\x*\x-1});
  \node[coulCourbe,font=\footnotesize] at (-1.2,3) {$\mathcal{C}_{1}$};
  \node[coulCourbeB,font=\footnotesize] at (2.9,2.6) {$\mathcal{C}_{2}$};
  \node[coulCourbeC,font=\footnotesize] at (2.6,-4.4) {$\mathcal{C}_{3}$};
\end{tikzpicture}
\end{minipage}
\end{methodebox}

\afaire{exercices 26 à 28}

\partiecours{Tableaux de variations}

\begin{proprietebox}[Propriété 8 --- Sommet et variations]
\begin{minipage}[c]{0.58\linewidth}\raggedright
Si $f$ a deux racines $x_{1}$ et $x_{2}$, le sommet $S$ de la parabole est \textbf{au
milieu} : son abscisse est
\[
  x_{S}=\dfrac{x_{1}+x_{2}}{2}
\]
et son ordonnée est $f(x_{S})$.
\end{minipage}\hfill
\begin{minipage}[c]{0.38\linewidth}\centering
\begin{tikzpicture}[x=0.42cm,y=0.42cm]
  \repere{-3}{-5}{5}{3}
  \gradx{-2,-1,1,2,3,4} \grady{-4,-2,2}
  \draw[coulCourbe,line width=1.1pt,domain=-1.9:3.9,samples=60] plot (\x,{(\x+1)*(\x-3)});
  \draw[red,dashed,thick] (1,-5)--(1,3);
  \path[red] (-1,0) node[croix]{} (3,0) node[croix]{};
  \node[red,above left,font=\footnotesize,inner sep=1pt] at (-1,0) {$x_{1}$};
  \node[red,above right,font=\footnotesize,inner sep=1pt] at (3,0) {$x_{2}$};
  \path (1,-4) node[croix]{} node[below right,font=\footnotesize,inner sep=2pt]{$S$};
\end{tikzpicture}
\end{minipage}

\smallskip
\begin{minipage}[t]{0.47\linewidth}\centering
{\small $a>0$ : $f(x_{S})$ est le \textbf{minimum}}\par\smallskip
\begin{tikzpicture}[scale=0.85]
  \tkzTabInit[lgt=1.6,espcl=2.2]{$x$/1,$f(x)$/1.6}{,$x_{S}$,}
  \tkzTabVar{+/ ,-/$f(x_{S})$,+/ }
\end{tikzpicture}
\end{minipage}\hfill
\begin{minipage}[t]{0.47\linewidth}\centering
{\small $a<0$ : $f(x_{S})$ est le \textbf{maximum}}\par\smallskip
\begin{tikzpicture}[scale=0.85]
  \tkzTabInit[lgt=1.6,espcl=2.2]{$x$/1,$f(x)$/1.6}{,$x_{S}$,}
  \tkzTabVar{-/ ,+/$f(x_{S})$,-/ }
\end{tikzpicture}
\end{minipage}
\end{proprietebox}

\begin{methodebox}[Méthode 10 --- Dresser le tableau de variations]
\etapeexemples
Donner le sommet de la parabole, puis le tableau de variations de :
\begin{questions}[label=\alph*)]
  \item $f(x)=2(x-1)(x+3)$
  \item $g(x)=-(x-2)(x-6)$
  \item $h(x)=3x^{2}-12x$
\end{questions}

\etape{À vous de jouer}
Donner le sommet de la parabole, puis le tableau de variations de :
\[
  f(x)=3(x-2)(x+4) \qquad g(x)=-2(x-1)(x-5) \qquad h(x)=-x^{2}+8x
\]

\end{methodebox}

\afaire{exercices 29 à 31}

\partiecours{Problèmes type épreuve anticipée}

À l'épreuve, une fonction du second degré sert de \textbf{modèle} (bénéfice, hauteur,
population…) : on calcule des images, on \textbf{vérifie} une forme factorisée donnée,
puis on en tire les racines, le sommet et les variations, et on \textbf{interprète}
chaque résultat dans le contexte.

\begin{methodebox}[Méthode 11 --- Étudier un modèle du second degré]
\etapeexemples
Un atelier fabrique et vend entre $0$ et $1\,200$ trottinettes électriques par an. On
note $x$ le nombre de trottinettes, en centaines ($x$ entre $0$ et $12$). Le
bénéfice, en milliers d'euros, est modélisé par
\[
  B(x)=-2x^{2}+24x-40 .
\]
Sa courbe est tracée ci-dessous ($x$ en centaines, $B(x)$ en milliers d'euros).
\begin{center}
\begin{tikzpicture}[x=0.7cm,y=0.4cm]
  \repere{0}{-4}{13}{5}
  \gradx{1,2,3,4,5,6,7,8,9,10,11,12}
  \grady{-4/-40,-3/-30,-2/-20,-1/-10,1/10,2/20,3/30,4/40}
  \draw[coulCourbe,line width=1.1pt,domain=0:12,samples=80] plot (\x,{(-2*\x*\x+24*\x-40)/10});
\end{tikzpicture}
\end{center}
\begin{questions}
  \item Par lecture graphique, déterminer pour quelles quantités le bénéfice est de
        $30\,000$~€.
  \item Calculer $B(1)$. Interpréter ce résultat dans le contexte de l'exercice.
  \item Vérifier que $B(x)=-2(x-2)(x-10)$.
  \item Pour quelles quantités le bénéfice est-il nul ?
  \item Dresser le tableau de variations de $B$ entre $0$ et $12$.
  \item Quel est le bénéfice maximal ? Pour combien de trottinettes ?
\end{questions}

\etape{À vous de jouer}
\begin{minipage}[c]{0.58\linewidth}\raggedright
On lance une fusée à eau depuis le sol. Sa hauteur, en mètres, au bout de $t$ secondes
($t$ entre $0$ et $4$) est
\[
  h(t)=-5t^{2}+20t .
\]
Sa courbe est tracée ci-contre ($t$ en secondes, $h(t)$ en mètres).
\end{minipage}\hfill
\begin{minipage}[c]{0.38\linewidth}\centering
\begin{tikzpicture}[x=1cm,y=0.45cm]
  \draw[grille,step=1] (0,0) grid (5,5);
  \draw[-{Stealth[length=2mm]},thick] (0,0)--(5.4,0) node[below left=-1pt]{$t$};
  \draw[-{Stealth[length=2mm]},thick] (0,0)--(0,5.4) node[below left=-1pt]{$h$};
  \node[below left,font=\scriptsize,inner sep=1.5pt] at (0,0) {$0$};
  \gradx{1,2,3,4}
  \grady{1/5,2/10,3/15,4/20}
  \draw[coulCourbe,line width=1.1pt,domain=0:4,samples=60] plot (\x,{(-5*\x*\x+20*\x)/5});
\end{tikzpicture}
\end{minipage}
\begin{questions}
  \item Par lecture graphique, déterminer à quels instants la fusée est à $15$~m de
        haut.
  \item Calculer $h(1)$. Interpréter ce résultat dans le contexte de l'exercice.
  \item Vérifier que $h(t)=-5t(t-4)$.
  \item À quels instants la fusée est-elle au sol ?
  \item Dresser le tableau de variations de $h$ entre $0$ et $4$.
  \item Quelle est la hauteur maximale ? À quel instant ?
\end{questions}

\end{methodebox}

\afaire{exercices 32 à 34}

\end{document}
